\section{Introduction}
\label{sec:intro}

% --- Background: ASR landscape ---
Automatic Speech Recognition (ASR) has progressed from hybrid HMM-DNN systems
\citep{hinton2012deep} to end-to-end sequence-to-sequence models such as
CTC \citep{graves2006ctc}, RNN-T \citep{graves2012rnnt}, and attention-based
encoder--decoders \citep{chan2016las}. Large-scale weakly-supervised models
such as Whisper \citep{radford2023whisper} further demonstrated that scaling
data and model size yields highly robust speech recognition across languages
and acoustic conditions.

% --- Motivation: why LLM-based ASR ---
In parallel, large language models (LLMs) have become the dominant paradigm for
text understanding and generation. A natural question is whether the strong
linguistic priors and in-context capabilities of LLMs can be leveraged to
improve ASR --- not only in raw word error rate (WER), but also in
\emph{contextual biasing}, \emph{multilingual transfer}, \emph{spoken language
understanding}, and \emph{instruction following}.

% --- This work ---
In this report we describe the design and training of \textbf{LLM-ASR}, a system
that connects a high-quality audio encoder to a pre-trained LLM via a
lightweight adapter, trained in three stages on a mixture of public and
in-house speech corpora. Our contributions are:

\begin{itemize}
  \item \textbf{Architecture.} A modular encoder--adapter--LLM design that
        keeps the LLM backbone frozen during early training and unlocks it
        progressively. (Section~\ref{sec:model})
  \item \textbf{Data.} A curated multilingual speech corpus covering
        \textit{N} hours across \textit{K} languages, with a transparent
        cleaning and deduplication pipeline. (Section~\ref{sec:data})
  \item \textbf{Training recipe.} A three-stage recipe (alignment,
        supervised fine-tuning, preference optimization) that scales
        stably from 1B to \textit{X}B parameters. (Section~\ref{sec:training})
  \item \textbf{Empirical results.} Strong results on standard ASR
        benchmarks and on contextual / instruction-driven variants, with
        ablations isolating the contribution of each component.
        (Sections~\ref{sec:experiments}--\ref{sec:analysis})
\end{itemize}

% --- Roadmap ---
The remainder of the report is organized as follows.
Section~\ref{sec:related} surveys related work.
Section~\ref{sec:data} describes the data pipeline.
Section~\ref{sec:model} introduces the architecture, and
Section~\ref{sec:training} details the training recipe.
Sections~\ref{sec:experiments} and \ref{sec:analysis} report results and
ablations. We discuss limitations and broader impact in
Section~\ref{sec:limitations} before concluding in
Section~\ref{sec:conclusion}.
